\documentclass[11pt]{article}
\usepackage[margin=1in]{geometry}
\usepackage{amsmath,amssymb,amsthm}
\usepackage{hyperref}

\newtheorem{theorem}{Theorem}
\newtheorem{remark}{Remark}

\newcommand{\E}{\mathbb E}
\newcommand{\Pp}{\mathbb P}
\newcommand{\R}{\mathbb R}

\title{A Jump Obstruction to Quadratic-Variation-Only Estimates}
\author{Run \texttt{fa\_banach\_001}, lane 7}
\date{2026-06-21}

\begin{document}
\maketitle

\section*{Source and Question}

The source paper is Mark Veraar and Ivan Yaroslavtsev,
\emph{Cylindrical continuous martingales and stochastic integration in
infinite dimensions}, arXiv:1602.03996.  At the end of the introduction the
authors ask what structure of a cylindrical \emph{noncontinuous} local
martingale is required in order to build a stochastic integration theory with
two-sided estimates for stochastic integrals.

In the continuous theory of the source paper, the main size datum is the
cylindrical quadratic variation \( [[M]] \), together with its operator
derivative.  The following elementary example shows that, after jumps are
allowed, a theory whose estimates see only the predictable second-order
quadratic characteristic cannot work.  Thus extra jump-size structure is not
cosmetic; it is forced already in the scalar case.

\section*{Idea of the Proof}

A rare jump of size \(R\) and intensity \(R^{-2}\) has unit predictable
quadratic variation.  Hence all these martingales look identical to any norm
that records only the second characteristic.  But for \(p>2\), the \(p\)-th
moment is dominated by the event that the rare jump occurs once, giving size
approximately \(R^p\) with probability approximately \(R^{-2}\).  The resulting
moment grows like \(R^{p-2}\).

\begin{theorem}
Let \(p>2\).  There is no constant \(C_p\) with the following property:
for every real-valued purely discontinuous square-integrable martingale \(M\)
on \([0,1]\) with predictable quadratic variation \(\langle M\rangle_t=t\),
\[
   \bigl(\E |M_1|^p\bigr)^{1/p} \leq C_p .
\]
Consequently, no noncontinuous cylindrical stochastic integration theory can
have uniform two-sided \(L^p\) estimates for the constant integrand if the
right-hand side depends only on the predictable quadratic variation/intensity
measure of the scalar projections.
\end{theorem}

\begin{proof}
Let \(R\geq 2\), and let \(N^R=(N^R_t)_{0\leq t\leq 1}\) be a Poisson process
with intensity \(\lambda_R=R^{-2}\).  Define
\[
   M^R_t = R\bigl(N^R_t-\lambda_R t\bigr), \qquad 0\leq t\leq 1.
\]
Then \(M^R\) is a real-valued square-integrable martingale with jumps of size
\(R\).  Since the predictable quadratic variation of the compensated Poisson
martingale \(N^R_t-\lambda_R t\) is \(\lambda_R t\), we have
\[
   \langle M^R\rangle_t = R^2\lambda_R t = t,
   \qquad 0\leq t\leq 1.
\]
Thus all martingales \(M^R\) have the same predictable second characteristic.

On the event \(\{N^R_1=1\}\),
\[
   |M^R_1| = R|1-\lambda_R| \geq R/2
\]
for \(R\geq 2\).  Also
\[
   \Pp(N^R_1=1)=\lambda_R e^{-\lambda_R}
      =R^{-2}e^{-R^{-2}} \geq e^{-1/4}R^{-2}.
\]
Therefore
\[
   \E |M^R_1|^p
   \geq (R/2)^p e^{-1/4}R^{-2}
   = 2^{-p}e^{-1/4} R^{p-2}.
\]
This tends to infinity as \(R\to\infty\), contradicting any bound depending
only on \(p\) and on the common value \(\langle M^R\rangle_1=1\).

The same example is a one-dimensional cylindrical martingale, so it embeds
unchanged into any broader cylindrical framework.  For the stochastic integral
of the constant scalar integrand \(1\), the integral is just \(M^R\), so the
failure is already a failure of the desired integral estimate.
\end{proof}

\begin{remark}[What the obstruction does not say]
This does not contradict Burkholder--Davis--Gundy inequalities for martingales
with jumps.  Those inequalities use richer data, for instance the optional
quadratic variation \([M^R]_1=R^2N^R_1\), whose \(L^{p/2}\)-size also grows
with \(R\).  Nor does it contradict jump integration theories with
Rosenthal-type or Poissonian norms.  The point is only that the continuous
paper's quadratic-variation-only philosophy cannot be transplanted to jumps
using only predictable second characteristics.
\end{remark}

\section*{Verification Notes}

The example is scalar, so no Banach-space geometry is hidden.  The predictable
quadratic variation computation is the standard one for compensated Poisson
processes.  The lower bound uses only the one-jump event and is uniform for
\(R\geq2\).

\section*{Novelty and Literature Check}

The run indexes were searched for arXiv:1602.03996, the source title,
``cylindrical noncontinuous'', ``two-sided estimates'', and ``scalar quadratic
variation''.  No existing run packet for this source was found.  A bounded web
search on 2026-06-21 for exact phrases around the source's \( [[M]] \)
question and tensor/scalar quadratic variation did not find a matching
packet-level answer.  The phenomenon is elementary and should be regarded as a
folklore-level obstruction unless a deeper novelty check says otherwise.

\section*{References}

\begin{itemize}
\item M. Veraar and I. Yaroslavtsev, \emph{Cylindrical continuous martingales
and stochastic integration in infinite dimensions}, Electron. J. Probab. 21
(2016), no. 59, 1--53; arXiv:1602.03996.
\item S. Dirksen, \emph{It\^o isomorphisms for \(L^p\)-valued Poisson
stochastic integrals}, Ann. Probab. 42 (2014), no. 6, 2595--2643.  Cited for
context on richer Poisson/Rosenthal jump norms, not used in the proof.
\end{itemize}

\section*{Human-Review Recommendation}

Verify the scope: this packet proves a negative result for
quadratic/intensity-only jump estimates.  It should not be read as a negative
answer to the full broad program of constructing noncontinuous cylindrical
integration theories with additional jump structure.

\end{document}
